\documentclass[11pt]{article}


\setlength{\parindent}{0pt}
\usepackage{xltxtra}
\usepackage{hyperref}
\hypersetup{hidelinks}
\usepackage{url}
\urlstyle{tt}
\usepackage{xcolor}
\definecolor{CVBlue}{RGB}{23,110,191}
\usepackage{calc}
\usepackage{graphicx}
\usepackage{tikz}
\usetikzlibrary{calc}
\usepackage{fontspec}
\usepackage{xeCJK}
\usepackage{enumitem}
\CJKsetecglue{} %% 取消中文与数字之间的间隙


%% 主文档字体设置
\setmainfont[
    Path = fonts/Main/,
    Extension = .otf,
    BoldFont = texgyretermes-bold.otf, % 加粗字体
]{texgyretermes-regular.otf} % 正文字体

% 中文字体设置
\setCJKmainfont[
    Path = fonts/hansans/,
    Extension = .ttf,
    BoldFont = NotoSansSC-Bold.ttf, % 加粗字体
]{NotoSansSC-Regular.otf} % 正文字体


\usepackage{relsize}
\usepackage{xspace}

% 使用 fontawesome（部分图标）
\usepackage{fontawesome} 

% A4纸，上下左右边距
\usepackage[
    a4paper,
    left=1.2cm,
    right=1.2cm,
    top=1.5cm,
    bottom=1cm,
    nohead
]{geometry}

\renewcommand{\baselinestretch}{1.5} % 行间距设为1.5

\usepackage{titlesec}
\usepackage{enumitem}
\setlist{noitemsep} % 取消列表项间的额外间距
%\setlist{nosep} % 取消所有垂直间距
\setlist[itemize]{topsep=0.25em, leftmargin=*}
\setlist[enumerate]{topsep=0.25em, leftmargin=*}

% --- 用于控制【不同项目之间】的垂直距离 ---
\newlength{\interProjectSpacing}
\setlength{\interProjectSpacing}{0.9em} % <--- 在此调整项目之间的距离
\newcommand{\projectsep}{\vspace{\interProjectSpacing}}

% --- 用于控制【项目标题】与下方【项目描述】的距离 ---
\newlength{\intraProjectTitleSep}
\setlength{\intraProjectTitleSep}{0.4em} % <--- 在此调整标题和描述的距离
\newcommand{\titlebreak}{\\[\intraProjectTitleSep]}

% --- 用于控制【项目描述】与下方【要点列表】的距离 ---
\newlength{\intraProjectListTopSep}
\setlength{\intraProjectListTopSep}{0.2em} % <--- 在此调整描述和列表的距离

% =======================================================================


\titleformat{\section}         % 定制 \section 命令 
{\large\bfseries\raggedright} % 将 section 标题设置为大号、粗体且左对齐
{}{0em}                      % 可用于为所有 section 添加前缀（如“章节...”）
{}                           % 可用于在标题前插入代码
[{\color{CVBlue}\titlerule}]  % 在标题后插入一条横线
\titlespacing*{\section}{0cm}{*1.6}{*1.2}



\begin{document}
\pagenumbering{gobble}

%%%% 利用tikz来定位照片
\begin{tikzpicture}[remember picture, overlay] 
    \node[anchor = north east] at ($(current page.north east)+(-2cm,-1.2cm)$) {\includegraphics[height=3cm]{avatar.jpg}};
  \end{tikzpicture}%
  %%%% 利用tikz来定位学校Logo，这里只在第一页显示
  \begin{tikzpicture}[remember picture, overlay] 
    \node[anchor = north west] at ($(current page.north west)+(0.5cm,+1.0cm)$) {\includegraphics[height=6cm]{zju.png}};
  \end{tikzpicture}%
\centerline{\LARGE\bfseries{樊奕铮}} 

\centerline{\normalsize{\faPhone\ 159-5755-1942 \quad \faEnvelopeO\ \href{mailto:3240103512@zju.edu.cn}{1928842988@qq.com}}} 


\section{\makebox[\widthof{\faGraduationCap}][c]{\color{CVBlue}\faGraduationCap}\ 教育背景}    
\textbf{浙江大学} \hfill 2024.9 -- 至今\\[0.5em] % 标题和正文间加一点距离
能源与环境系统工程\quad 大二 
\begin{itemize}[nosep]
    \item 相关课程：《工程流体力学》、《工程热力学》、《c程序设计基础及实验》、《电工电子学》、《工程力学》
\end{itemize}

\section{\makebox[\widthof{\faUsers}][c]{\color{CVBlue}\faUsers}\ 项目经历}

% --- 第一个项目 ---
% 将标题行末尾的 \\ 替换为 \titlebreak 命令
\textbf{AI+能源智能关灯机器人} \hfill 2026.05 -- 2026.07 \titlebreak
项目描述：面向校园场景用电浪费问题，设计并实现一款基于 AI 视觉识别的智能关灯节能机器人，通过人体检测与环境感知自动管控照明设备，实现无人关灯、智能节电，有效降低场景能耗。
% 在 itemize 的选项中，使用 topsep=\intraProjectListTopSep 来控制上边距
\begin{itemize}[nosep, topsep=\intraProjectListTopSep]
    \item \textbf{AI模型训练}：基于小车自带的yolo训练人体检测，灯光检测，开关检测模型。
    \item \textbf{代码编写}：编写机器人路网规划，智能关灯代码，构建环境。
\end{itemize}
\bigskip
\bigskip
\textbf{ai管理碳中和的信任度危机} \hfill 2025.08 -- 2026.02 \titlebreak
项目描述：在 “双碳” 目标背景下，AI 技术在碳排放核算、碳交易监管等场景中应用日益广泛，但数据透明度、算法黑箱等问题引发了社会对碳中和治理信任度的担忧。本项目旨在分析 AI 管控碳中和过程中的信任度危机成因，构建一套可量化的信任度评估框架，为提升碳中和治理的公信力提供决策参考。
% 在 itemize 的选项中，使用 topsep=\intraProjectListTopSep 来控制上边距
\begin{itemize}[nosep, topsep=\intraProjectListTopSep]
    \item \textbf{文献调研}：系统梳理国内外关于 AI 治理、碳核算、信任机制的研究现状，整理形成文献综述。
    \item \textbf{数据收集与分析}：收集典型案例中的碳排放数据、AI 决策记录，使用 Python 进行数据清洗与统计分析。
\end{itemize}
\bigskip
\textbf{基于机器学习算法的热负荷预测，电负荷预测、环境温度预测研究} \hfill 2026.03 -- 至今 \titlebreak
项目描述：本项目面向综合能源系统用能预测需求，依托机器学习算法开展电负荷、热负荷与环境温度联合预测研究。项目完成数据预处理、特征筛选、多类预测模型搭建与精度对比实验，挖掘气象、时序与冷热负荷间耦合规律，构建高精度多任务预测模型，为园区能源调度、节能管控提供数据支撑。
% 在 itemize 的选项中，使用 topsep=\intraProjectListTopSep 来控制上边距
\begin{itemize}[nosep, topsep=\intraProjectListTopSep]
    \item \textbf{机器学习预测模型搭建与训练}：基于Python搭建多类模型，分别完成传统机器学习、时序深度学习、多任务学习框架和网格搜索。
    \item \textbf{模型对比实验与结果分析}：分别单独预测环境温度、电负荷、热负荷，记录各模型误差指标；对比单任务预测与多任务联合预测精度差异；绘制负荷真实值 - 预测值对比曲线、误差柱状图、特征重要度分布图；分析温度变化对电、热负荷波动的影响规律，解释不同模型适用场景。
\end{itemize}



\section{\makebox[\widthof{\faCogs}][c]{\color{CVBlue}\faCogs}\ 技术栈}
\begin{itemize}[nosep]
    \item \textbf{编程语言：} JavaScript, Python, C++
    \item \textbf{开发工具：} Git, Vim,LaTex，claude code
\end{itemize}
\section{\makebox[\widthof{\faGraduationCap}][c]{\color{CVBlue}\faList}\ 获奖情况}
\begin{itemize}
    \item 大学生物理创新竞赛三等奖
    \item AI+能源大学生创新竞赛二等奖
\end{itemize}
    
\section{\makebox[\widthof{\faInfo}][c]{\color{CVBlue}\faInfo}\ 其他}
\begin{itemize}[parsep=0.5ex]
    \item \textbf{英语水平：} CET-4, CET-6
    \item 基础人工智能技术，有效使用claudecode辅助作业
    \item steam多个热门mod作者
\end{itemize}
\end{document}
